% from ega2-1-fr.tex
\oldpage[II]{89}

et on conclut de (\hyperref[II.4.5.13.2-fr]{4.5.13.2}) et
(\hyperref[II.4.5.13.4-fr]{4.5.13.4}) que l'on a
$\partial(g^{\otimes(m+1)})=0$~; $g^{\otimes(m+1)}$ est donc l'image
canonique d'une section $f$ de $\mathcal{L}^{\otimes n(m+1)}$ au-dessus de
$X$, ce qui achève de démontrer (\hyperref[II.4.5.13-fr]{4.5.13}).
\end{proof}

\begin{corollary}[4.5.14]
\label{II.4.5.14-fr}
Soient $X$ un schéma noethérien, $j$ l'injection canonique
$X_{\mathrm{red}}\to X$. Pour qu'un $\mathcal{O}_X$-Module inversible
$\mathcal{L}$ soit ample, il faut et il suffit que $j^*(\mathcal{L})$ soit
un $\mathcal{O}_{X_{\mathrm{red}}}$-Module ample.
\end{corollary}

Cela résulte de (\hyperref[I.6.1.6-fr]{I, 6.1.6}).

\subsection{Faisceaux relativement amples}
\label{II.4.6-fr}

\begin{definition}[4.6.1]
\label{II.4.6.1-fr}
Soient $f:X\to Y$ un morphisme quasi-compact, $\mathcal{L}$ un
$\mathcal{O}_X$-Module inversible. On dit que $\mathcal{L}$ est ample
relativement à $f$, ou relativement à $Y$, ou $f$-ample, ou $Y$-ample (ou
même ample si aucune confusion n'en résulte avec la notion définie dans
(\hyperref[II.4.5.3-fr]{4.5.3})) s'il existe un recouvrement $(U_\alpha)$
de $Y$ formé d'ouverts affines tel que si on pose
$X_\alpha=f^{-1}(U_\alpha)$, $\mathcal{L}|X_\alpha$ soit un
$\mathcal{O}_{X_\alpha}$-Module ample pour tout $\alpha$.
\end{definition}

L'existence d'un $\mathcal{O}_X$-Module $f$-ample entraîne que $f$ est
nécessairement séparé ((\hyperref[II.4.5.3-fr]{4.5.3}) et
(\hyperref[I.5.5.5-fr]{I, 5.5.5})).

\begin{proposition}[4.6.2]
\label{II.4.6.2-fr}
Soit $f:X\to Y$ un morphisme quasi-compact, $\mathcal{L}$ un
$\mathcal{O}_X$-Module inversible. Si $\mathcal{L}$ est très ample
relativement à $f$, il est ample relativement à $f$.
\end{proposition}

Cela résulte du caractère local (sur $Y$) de la notion de faisceau très
ample (\hyperref[II.4.4.5-fr]{4.4.5}), de la définition
(\hyperref[II.4.6.1-fr]{4.6.1}) et du critère
(\hyperref[II.4.5.10-fr]{4.5.10, (i)}).

\begin{proposition}[4.6.3]
\label{II.4.6.3-fr}
Soient $f:X\to Y$ un morphisme quasi-compact, $\mathcal{L}$ un
$\mathcal{O}_X$-Module inversible, et soit $\mathcal{S}$ la
$\mathcal{O}_Y$-Algèbre graduée
$\bigoplus_{n\geq0}f_*(\mathcal{L}^{\otimes n})$. Les conditions suivantes
sont équivalentes~:
\begin{enumerate}
  \item[\rm{a)}] $\mathcal{L}$ est $f$-ample.
  \item[\rm{b)}] $\mathcal{S}$ est quasi-cohérente et l'homomorphisme
  canonique $\sigma:f^*(\mathcal{S})\to
  \bigoplus_{n\geq0}\mathcal{L}^{\otimes n}$
  (\hyperref[0.4.4.3-fr]{0, 4.4.3}) est tel que le $Y$-morphisme
  $r_{\mathcal{L},\sigma}:G(\sigma)\to\operatorname{Proj}(\mathcal{S})=P$
  soit partout défini et soit une immersion ouverte dominante.
  \item[\rm{b')}] Le morphisme $f$ est séparé, le $Y$-morphisme
  $r_{\mathcal{L},\sigma}$ est partout défini et est un homéomorphisme de
  l'espace sous-jacent à $X$ sur un sous-espace de
  $\operatorname{Proj}(\mathcal{S})$.
\end{enumerate}
En outre, lorsqu'il en est ainsi, pour tout $n\in\mathbf{Z}$,
l'homomorphisme canonique
\begin{equation}
\label{II.4.6.3.1-fr}
  r_{\mathcal{L},\sigma}^*(\mathcal{O}_P(n))
  \longrightarrow\mathcal{L}^{\otimes n}
\tag{4.6.3.1}
\end{equation}
défini dans (\hyperref[II.3.7.9.1-fr]{3.7.9.1}), est un isomorphisme.

Enfin, pour tout $\mathcal{O}_X$-Module quasi-cohérent $\mathcal{F}$, si
l'on pose
$\mathcal{M}=\bigoplus_{n\geq0}f_*(\mathcal{F}\otimes
\mathcal{L}^{\otimes n})$, l'homomorphisme canonique
\begin{equation}
\label{II.4.6.3.2-fr}
  r_{\mathcal{L},\sigma}^*(\widetilde{\mathcal{M}})
  \longrightarrow\mathcal{F}
\tag{4.6.3.2}
\end{equation}
défini dans (\hyperref[II.3.7.9.2-fr]{3.7.9.2}), est un isomorphisme.
\end{proposition}

On a remarqué que a) implique que $f$ est séparé, donc $\mathcal{S}$ est
quasi-cohérente (\hyperref[I.9.2.2-fr]{I, 9.2.2, a)}). Comme le fait que
$r_{\mathcal{L},\sigma}$ soit une immersion partout définie est de
caractère local sur $Y$, pour prouver que a) entraîne b), on peut supposer
$Y$ affine et $\mathcal{L}$ ample~;

